% !TEX root = ../main.tex
\section{Considerações Pessoais e Agradecimentos}
\label{sec:consideracoes_pessoais}

A participação no projeto \textbf{Bio-Waste2Carbon} no âmbito do concurso \textbf{Capture the Future II} representou uma oportunidade ímpar de desenvolvimento académico, técnico e pessoal. Ao longo destes três meses, foi-me possível aplicar e consolidar, num ambiente de investigação aplicada e interdisciplinar de elevado nível, os conhecimentos teóricos adquiridos durante o meu percurso formativo em Engenharia Mecânica, articulando desafios de sustentabilidade ambiental, automação, visão computacional e transição energética.

A complementaridade entre a experiência prática de bancada nos laboratórios da ADAI e da Universidade de Coimbra (durante o mês de julho) e o desenvolvimento computacional avançado em regime remoto (durante os meses de agosto e setembro) proporcionou-me uma perspetiva holística da atividade de engenharia e investigação. A superação de dificuldades técnicas reais --- desde a modelação tridimensional em CAD e integração mecânica de subsistemas na bancada experimental até à resolução de desafios complexos de inteligência artificial em Python, cinemática de alta velocidade e modelação física tridimensional --- reforçou substancialmente a minha autonomia, rigor metodológico e capacidade de pensamento crítico.

No âmbito dos agradecimentos, manifesto a minha mais profunda e sincera gratidão ao meu orientador, o \textbf{Eng.~Rodrigo Branco}, pelo acompanhamento permanente, pela partilha inexaurível de conhecimentos, pela exigência e rigor científico com que sempre norteou os trabalhos, e pelo constante estímulo e confiança depositados em mim desde o primeiro dia de estágio. O ambiente de colaboração, proximidade e incentivo proporcionado foi um fator determinante para a concretização de todos os objetivos traçados.

Agradeço igualmente a toda a equipa de investigação do projeto Bio-Waste2Carbon, à \textbf{ADAI} (Associação para o Desenvolvimento da Aerodinâmica Industrial), à empresa líder \textbf{SustainC}, à \textbf{Fravizel} e ao Departamento de Engenharia Mecânica da \textbf{Universidade de Coimbra}, por disponibilizarem os recursos técnicos e laboratoriais que tornaram possível a realização deste trabalho. Um agradecimento extensivo aos colegas e investigadores dos laboratórios de energia pela receção exemplar e espírito de entreajuda demonstrado durante a fase experimental.

Concluo este estágio com um sentimento profundo de realização, gratidão e dever cumprido, consciente de que esta experiência constitui uma base sólida e estruturante para o meu futuro percurso académico e profissional na engenharia, motivando-me a continuar a dedicar o meu esforço à inovação tecnológica e à investigação científica de excelência em Portugal.
